\documentclass[10pt,a4paper]{amsart}
\usepackage[margin=19mm]{geometry}
\usepackage{amsmath,amssymb,mathtools}
\usepackage[scr=rsfs,cal=euler]{mathalfa}
\usepackage{xfrac}
\usepackage[unicode,hidelinks,bookmarks=false]{hyperref}
\newcommand{\ie}{i.e.\ }
\newcommand{\cf}{cf.\ }
\newcommand{\resp}{respectively\ }
\newcommand{\Subsection}{\subsection}
\providecommand{\nobreakdash}{\nobreak\hbox{-}}
\setlength{\emergencystretch}{2em}
\multlinegap0pt
\usepackage{amssymb}
\usepackage{dsfont}
\usepackage{bm}
\usepackage{enumerate}
\newtheorem{theorem}{Theorem}
\title{$RLL$-REALIZATION OF TWO-PARAMETER QUANTUM AFFINE ALGEBRA IN TYPE $C_n^{(1)}$}
\author{Xin Zhong, Naihong Hu, Naihuan Jing$^*$}
\date{Source: arXiv:2405.06597v2 (CC BY 4.0)}
\begin{document}
\maketitle

\noindent\emph{Source and licence.} $RLL$-REALIZATION OF TWO-PARAMETER QUANTUM AFFINE ALGEBRA IN TYPE $C_n^{(1)}$. Version arXiv:2405.06597v2 (CC BY 4.0), distributed by arXiv under the Creative Commons Attribution 4.0 International licence, http://creativecommons.org/licenses/by/4.0/, as recorded in the arXiv licence metadata for this exact version and shown on the abstract page https://arxiv.org/abs/2405.06597v2. Full licence notice: \texttt{licenses/arxiv-2405.06597-CC-BY-4.0.txt}.\par\smallskip

\noindent\textit{Editorial scope.} Counted unit: the article's theorem theorem (source0.tex lines 264--290). The article's complete original proof of it is delivered with it (source0.tex lines 291--387, one \texttt{proof} environment of 97 lines). No mathematics is written, repaired or re-derived: the statement and the proof are the article's own lines, in the article's own notation, and no retained line is substituted or re-flowed.  No further source-local context is needed: every cross-reference of every retained line resolves inside the two delivered spans.  Nothing was omitted from any retained span, so the package carries no omission record.  No retained line contains a citation, so the wrapper carries no bibliography and no bibliography entry.  No retained line contains a figure environment, an image-inclusion command or drawing code, so no image is delivered and the package declares no asset dependency.  Editorial formatting only, in addition: an AMS \texttt{amsart} preamble that restates the article's own declarations and macros used by the retained lines, namely the environments \texttt{theorem} on separate counters, together with the standard packages those lines need.  The retained environments print in this document as Theorem 1 in order of appearance, whereas the article prints the same items with the numbering of its own full document.  The complete native source of the article as distributed in the arXiv e-print for this exact version is bundled unchanged as \texttt{sources/158979/source0.tex}.
			\begin{theorem} The $R$-matrix associated with $T_1$ is given by
	\begin{flalign*}
					& R=r^{\frac{1}{2}} s^{-\frac{1}{2}} \sum_{i=1}^{2 n} E_{i i} \otimes E_{i i}+r^{-\frac{1}{2}} s^{\frac{1}{2}} \sum_{i=1}^{2 n} E_{i i} \otimes E_{i^{'} i^{'}}+r^{\frac{1}{2}} s^{\frac{1}{2}}\left\{\sum_{\substack{1 \leq i \leq n-1 \\
							i+1 \leq j \leq n}} E_{i i} \otimes E_{j j}\right. \\
					& +\sum_{\substack{1 \leq i \leq n-1 \\
							i^{'}+1 \leq j \leq 2 n}} E_{i i} \otimes E_{j j}+\sum_{j=n+2}^{2 n} E_{n n} \otimes E_{j j}+\sum_{\substack{n+1 \leq j \leq 2 n-1 \\
							j+1 \leq i \leq 2 n}} E_{i i} \otimes E_{j j} \\
					& \left.+\sum_{\substack{1 \leq i \leq n-1 \\
							n+1 \leq j \leq 2 n-i}} E_{j j} \otimes E_{i i}\right\}+r^{-\frac{1}{2}} s^{-\frac{1}{2}}\left\{\sum_{\substack{1 \leq i \leq n-1 \\
							i+1 \leq j \leq n}} E_{j j} \otimes E_{i i}+\sum_{\substack{1 \leq i \leq n-1 \\
							i^{'}+1 \leq j \leq 2 n}} E_{j j} \otimes E_{i i}\right.
	\end{flalign*}
	\begin{flalign*}
	& \left.+\sum_{j=n+2}^{2 n} E_{j j} \otimes E_{n n}+\sum_{\substack{n+1 \leq j \leq 2 n-1 \\
			j+1 \leq i \leq 2 n}} E_{j j} \otimes E_{i i}+\sum_{\substack{1 \leq i \leq n-1 \\
			n+1 \leq j \leq 2 n-i}} E_{i i} \otimes E_{j j}\right\} \\
	& +\left(r^{\frac{1}{2}} s^{-\frac{1}{2}}-r^{-\frac{1}{2}} s^{\frac{1}{2}}\right)\left\{\sum_{i<j} E_{i j} \otimes E_{j i}-\sum_{i>j}\left(r^{\frac{1}{2}} s^{-\frac{1}{2}}\right)^{\bar{\imath}-\bar{\jmath}} \varepsilon_i \varepsilon_j E_{i^{\prime} j^{\prime}} \otimes E_{i j}\right\},
	\end{flalign*}
where
$$
\varepsilon_i=\left\{\begin{aligned}
	1, & \quad\text { for } \quad i=1, \ldots, n, \\
	-1, & \quad\text { for } \quad i=n+1, \ldots, 2 n,
\end{aligned}\right.
$$
and $(\overline{1}, \overline{2}, \ldots, \overline{2 n})=(n, n-1, \ldots, 1,-1, \ldots,-n)$.
\end{theorem}

\begin{proof}
	For $i=1,2, \cdots ,n-1$, we have $\Delta\left(e_i\right)=e_i \otimes \omega_i+1 \otimes e_i$ with the matrix form given by
\begin{equation}\label{e:Delta-ei}
\begin{aligned}
	&r^{-\frac{3}{2}} s^{-\frac{1}{2}} E_{i, i} \otimes E_{i+1, i}+r^{-\frac{1}{2}} s^{-\frac{3}{2}} E_{i+1, i+1} \otimes E_{i+1, i}+r^{-\frac{1}{2}} s^{\frac{1}{2}} E_{2 n-i, 2 n-i} \otimes E_{i+1, i}\\
	&+r^{\frac{1}{2}} s^{-\frac{1}{2}} E_{2 n-i+1,2 n-i+1} \otimes E_{i+1, i}+r^{-\frac{1}{2}} s^{-\frac{1}{2}} \sum_{j \neq\{i, i+1,2 n-i, 2 n-i+1\}} E_{j, j} \otimes E_{i+1, i}\\
	&-r^{-1} E_{i, i} \otimes E_{2 n-i+1,2 n-i}-s^{-1} E_{i+1, i+1} \otimes E_{2 n-i+1,2 n-i}-s E_{2 n-i, 2 n-i} \otimes E_{2 n-i+1,2 n-i}\\
	&-r E_{2 n-i+1,2 n-i+1} \otimes E_{2 n-i+1,2 n-i}-\sum_{j \neq\{i, i+1,2 n-i, 2 n-i+1\}} E_{j, j} \otimes E_{2 n-i+1,2 n-i}\\
	&+r^{-\frac{1}{2}} s^{-\frac{1}{2}} E_{i+1, i} \otimes\left(\sum_{j=1}^{2 n} E_{j, j}\right)-E_{2 n-i+1,2 n-i} \otimes\left(\sum_{j=1}^{2 n} E_{j, j}\right).
\end{aligned}
\end{equation}	
and $\Delta^{c o p}\left(e_i\right)=\omega_i \otimes e_i+e_i \otimes 1$, whose matrix is the flip of \eqref{e:Delta-ei}.
Here the flip means $A\otimes B\mapsto B\otimes A$ term by term.
%$$
%r^{-\frac{3}{2}} s^{-\frac{1}{2}} E_{i+1, i} \otimes E_{i, i}+r^{-\frac{1}{2}} s^{-\frac{3}{2}} E_{i+1, i} \otimes E_{i+1, i+1}+r^{-\frac{1}{2}} s^{\frac{1}{2}} E_{i+1, i} \otimes E_{2 n-i, 2 n-i}
%$$
%$$
%+r^{\frac{1}{2}} s^{-\frac{1}{2}} E_{i+1, i} \otimes E_{2 n-i+1,2 n-i+1}+r^{-\frac{1}{2}} s^{-\frac{1}{2}} E_{i+1, i} \otimes \sum_{j \neq\{i, i+1,2 n-i, 2 n-i+1\}} E_{j, j}
%$$
%$$
%-r^{-1} E_{2 n-i+1,2 n-i} \otimes E_{i, i}-s^{-1} E_{2 n-i+1,2 n-i} \otimes E_{i+1, i+1}-s E_{2 n-i+1,2 n-i} \otimes E_{2 n-i, 2 n-i}
%$$
%$$
%-r E_{2 n-i+1,2 n-i} \otimes E_{2 n-i+1,2 n-i+1}-E_{2 n-i+1,2 n-i} \otimes \sum_{j \neq\{i, i+1,2 n-i, 2 n-i+1\}} E_{j, j}
%$$
%$$
%+r^{-\frac{1}{2}} s^{-\frac{1}{2}}\left(\sum_{j=1}^{2 n} E_{j, j}\right) \otimes E_{i+1, i}-\left(\sum_{j=1}^{2 n} E_{j, j}\right) \otimes E_{2 n-i+1,2 n-i}.
%$$
Then we can check that
$\Delta^{c o p}\left(e_i\right) R=R \Delta\left(e_i\right)$, for $1\leq i\leq n-1$. Here we only consider nontrivial equations.

To see this, note that for $i=1,2, \cdots ,n-1$, one has that
$$
\left(r^{-\frac{1}{2}} s^{-\frac{3}{2}} E_{i+1, i} \otimes E_{i+1, i+1}\right)\left(r^{\frac{1}{2}} s^{-\frac{1}{2}}-r^{-\frac{1}{2}} s^{\frac{1}{2}}\right)\left(E_{i, i+1} \otimes E_{i+1, i}\right)
$$
$$
+\left(r^{-\frac{1}{2}} s^{-\frac{1}{2}} E_{i+1, i+1} \otimes E_{i+1, i}\right)\left(r^{-\frac{1}{2}} s^{-\frac{1}{2}} E_{i+1, i+1} \otimes E_{i, i}\right)
$$
$$
=\left(r^{\frac{1}{2}} s^{-\frac{1}{2}} E_{i+1, i+1} \otimes E_{i+1, i+1}\right)\left(r^{-\frac{1}{2}} s^{-\frac{3}{2}} E_{i+1, i+1} \otimes E_{i+1, i}\right),
$$

Also we compute that
$$
\left(r^{\frac{1}{2}} s^{-\frac{1}{2}} E_{i+1, i} \otimes E_{2 n-i+1,2 n-i+1}\right)\left(r^{\frac{1}{2}} s^{-\frac{1}{2}}-r^{-\frac{1}{2}} s^{\frac{1}{2}}\right)\left[\left(1+\left(r^{\frac{1}{2}} s^{-\frac{1}{2}}\right)^{\rho_1}\right) E_{i, 2 n-i+1} \otimes E_{2 n-i+1, i}\right]
$$
$$
+\left(-E_{i+1, i+1} \otimes E_{2 n-i+1,2 n-i}\right)\left(r^{\frac{1}{2}} s^{-\frac{1}{2}}-r^{-\frac{1}{2}} s^{\frac{1}{2}}\right)\left[\left(r^{\frac{1}{2}} s^{-\frac{1}{2}}\right)^{\rho_2} E_{i+1,2 n-i+1} \otimes E_{2 n-i, i}\right]
$$
$$
=\left[\left(r^{\frac{1}{2}} s^{-\frac{1}{2}}-r^{-\frac{1}{2}} s^{\frac{1}{2}}\right) E_{i+1,2 n-i+1} \otimes E_{2 n-i+1, i+1}\right]\left(r^{\frac{1}{2}} s^{-\frac{1}{2}} E_{2 n-i+1,2 n-i+1} \otimes E_{i+1, i}\right),
$$
where $
\rho_1=\overline{2 n-i+1}-\bar{i}$ and $\rho_2=\overline{2 n-i}-\bar{i}.$\\

Next, for $i=1,2, \cdots ,n-1, j \neq 2 n-i, 2 n-i+1$, and $j>i+1,$
$$
\left(r^{-\frac{1}{2}} s^{-\frac{1}{2}} E_{i+1, i} \otimes E_{j, j}\right)\left[\left(r^{\frac{1}{2}} s^{-\frac{1}{2}}-r^{-\frac{1}{2}} s^{\frac{1}{2}}\right) E_{i, j} \otimes E_{j, i}\right]
$$
$$
=\left[\left(r^{\frac{1}{2}} s^{-\frac{1}{2}}-r^{-\frac{1}{2}} s^{\frac{1}{2}}\right) E_{i+1, j} \otimes E_{j, i+1}\right]\left(r^{-\frac{1}{2}} s^{-\frac{1}{2}} E_{j, j} \otimes E_{i+1, i}\right),
$$
and for $i=1,2, \cdots ,n-2$,
$$
\left(-E_{i+1, i+1} \otimes E_{2 n-i+1,2 n-i}\right)\left(r^{-\frac{1}{2}} s^{\frac{1}{2}} E_{i+1, i+1} \otimes E_{2 n-i, 2 n-i}\right)
$$
$$
+\left(r^{\frac{1}{2}} s^{-\frac{1}{2}} E_{i+1, i} \otimes E_{2 n-i+1,2 n-i+1}\right)\left[-\left(r^{\frac{1}{2}} s^{-\frac{1}{2}}-r^{-\frac{1}{2}} s^{\frac{1}{2}}\right)\left(r^{\frac{1}{2}} s^{-\frac{1}{2}}\right)^{\rho_3} E_{i, i+1} \otimes E_{2 n-i+1,2 n-i}\right]
$$
$$
=\left(r^{\frac{1}{2}} s^{\frac{1}{2}} E_{i+1, i+1} \otimes E_{2 n-i+1,2 n-i+1}\right)\left(-s^{-1} E_{i+1, i+1} \otimes E_{2 n-i+1,2 n-i}\right),
$$
where $\rho_3=\overline{2 n-i+1}-\overline{2 n-i}.$\\

 Finally, we can check that for $i=1,2, \cdots ,n-1,$
$$
\left(-E_{i+1, i+1} \otimes E_{2 n-i+1,2 n-i}\right)\left(r^{\frac{1}{2}} s^{-\frac{1}{2}}-r^{-\frac{1}{2}} s^{\frac{1}{2}}\right)\left[\left(1+\left(r^{\frac{1}{2}} s^{-\frac{1}{2}}\right)^{\rho_4}\right) E_{i+1,2 n-i} \otimes E_{2 n-i, i+1}\right]
$$
$$
+\left(r^{\frac{1}{2}} s^{-\frac{1}{2}} E_{i+1, i} \otimes E_{2 n-i+1,2 n-i+1}\right)\left(r^{\frac{1}{2}} s^{-\frac{1}{2}}-r^{-\frac{1}{2}} s^{\frac{1}{2}}\right)\left[\left(r^{\frac{1}{2}} s^{-\frac{1}{2}}\right)^{\rho_5} E_{i, 2 n-i} \otimes E_{2 n-i+1, i+1}\right]
$$
$$
=\left[\left(r^{\frac{1}{2}} s^{-\frac{1}{2}}-r^{-\frac{1}{2}} s^{\frac{1}{2}}\right) E_{i+1,2 n-i+1} \otimes E_{2 n-i+1, i+1}\right]\left(-E_{2 n-i+1,2 n-i} \otimes E_{i+1, i+1}\right),
$$
and also
$$
\left(-s^{-1} E_{2 n-i+1,2 n-i} \otimes E_{i+1, i+1}\right)\left(r^{-\frac{1}{2}} s^{\frac{1}{2}} E_{2 n-i, 2 n-i} \otimes E_{i+1, i+1}\right)
$$
$$
=\left(r^{-\frac{1}{2}} s^{-\frac{1}{2}} E_{2 n-i+1,2 n-i+1} \otimes E_{i+1, i+1}\right)\left(-E_{2 n-i+1,2 n-i} \otimes E_{i+1, i+1}\right).
$$
These relations  imply that the $R$-matrix does satisfy $\Delta^{c o p}\left(e_i\right) R=R \Delta\left(e_i\right)$,for $1\leq i\leq n-1$.
Similarly we can check that $
\Delta^{c o p}\left(e_n\right) R=R \Delta\left(e_n\right)$ holds.  Note that the corresponding equations for $f_i$ immediately follow
by $\tau$ and the flip operation on tensor product. Finally, the equations $\Delta^{c o p}\left(w_i\right) R=R \Delta\left(w_i\right)$ and $\Delta^{c o p}\left(w_i^{\prime}\right) R=R \Delta\left(w_i^{\prime}\right)$ are trivial, thus we have shown that the given $R$-matrix is correct.
%for $i=1,2 \cdots n,$ which is trivial.
\end{proof}

\end{document}
